\begin{figure*}[t]
  \centering
  \resizebox{\linewidth}{!}{%
  \begin{tikzpicture}[
    >=Latex,
    font=\small,
    paneltitle/.style={font=\bfseries},
    treelabel/.style={font=\bfseries\small},
    internal/.style={
      circle,
      draw=blue!55!black,
      line width=0.55pt,
      fill=blue!5,
      minimum size=10.2mm,
      inner sep=1pt,
      align=center
    },
    leaf/.style={
      rounded corners=2pt,
      draw=green!50!black,
      line width=0.55pt,
      fill=green!10,
      minimum width=14mm,
      minimum height=8.6mm,
      inner sep=2pt,
      align=center
    },
    note/.style={align=center, font=\scriptsize},
    bigtable/.style={
      matrix of nodes,
      ampersand replacement=\&,
      row sep=0.75mm,
      column sep=\colgap,
      nodes={
        anchor=center,
        inner xsep=2.5pt,
        inner ysep=1.2pt,
        minimum height=4.5mm,
        text depth=.25ex,
        outer sep=0pt
      },
      row 1/.style={nodes={font=\bfseries}},
      column 1/.style={nodes={minimum width=7mm}},
      column 2/.style={nodes={minimum width=11mm}},
      column 3/.style={nodes={minimum width=10mm}},
      column 4/.style={nodes={minimum width=13mm}},
      column 5/.style={nodes={minimum width=9mm}},
      column 6/.style={nodes={minimum width=13mm}},
      column 7/.style={nodes={minimum width=13mm}},
      column 8/.style={nodes={minimum width=12mm}}
    },
    colbg/.style={
      draw=gray!70,
      rounded corners=2pt,
      line width=.5pt,
      fill=white,
      inner sep=2.8pt
    },
    meta/.style={
      draw=gray!70,
      rounded corners=2pt,
      fill=gray!6,
      align=center,
      inner sep=6pt
    },
    fetchrow/.style={
      draw=red!75!black,
      rounded corners=2pt,
      line width=.85pt,
      inner xsep=4pt,
      inner ysep=2pt
    },
    fetchtxt/.style={text=red!75!black, font=\bfseries},
    outer/.style={
      draw=gray!55,
      dashed,
      rounded corners=4pt,
      inner sep=8pt
    }
  ]

    \def\colgap{4.2mm}

    % =========================
    % Left panel: tree semantics
    % =========================
    \node[paneltitle] (ltitle) at (-4.8,3.12) {(a) Tree semantics};

    % Tree 1
    \node[treelabel] at (-7.55,2.42) {tree $1$};
    \node[internal] (t0r) at (-7.55,1.72) {$1$\\[-1mm]\scriptsize$x_1<3$};
    \node[internal] (t0l) at (-9.05,0.38) {$2$\\[-1mm]\scriptsize$x_2<4$};
    \node[leaf]     (t0rleaf) at (-6.00,0.38) {$7$\\[-1mm]\scriptsize$w_2$};
    \node[leaf]     (t0w0) at (-10.00,-1.02) {$5$\\[-1mm]\scriptsize$w_0$};
    \node[leaf]     (t0w1) at (-8.35,-1.02) {$6$\\[-1mm]\scriptsize$w_1$};

    \draw[->, line width=.75pt] (t0r) -- (t0l);
    \draw[->, line width=.75pt] (t0r) -- (t0rleaf);
    \draw[->, line width=.75pt] (t0l) -- (t0w0);
    \draw[->, line width=.75pt] (t0l) -- (t0w1);

    % Tree 2
    \node[treelabel] at (-2.65,2.42) {tree $2$};
    \node[internal] (t1r) at (-2.65,1.72) {$3$\\[-1mm]\scriptsize$x_1<2$};
    \node[leaf]     (t1lleaf) at (-4.15,0.38) {$8$\\[-1mm]\scriptsize$w'_0$};
    \node[internal] (t1rr) at (-1.15,0.38) {$4$\\[-1mm]\scriptsize$x_2<6$};
    \node[leaf]     (t1w1) at (-2.05,-1.02) {$9$\\[-1mm]\scriptsize$w'_1$};
    \node[leaf]     (t1w2) at (-0.40,-1.02) {$10$\\[-1mm]\scriptsize$w'_2$};

    \draw[->, line width=.75pt] (t1r) -- (t1lleaf);
    \draw[->, line width=.75pt] (t1r) -- (t1rr);
    \draw[->, line width=.75pt] (t1rr) -- (t1w1);
    \draw[->, line width=.75pt] (t1rr) -- (t1w2);

    % =====================================
    % Right panel: one unified committed table
    % =====================================
    \node[paneltitle] (rtitle) at (7.25,3.12) {(b) Committed row view used by the proof};

    \matrix[bigtable, anchor=north west] (tbl) at (1.15,2.55) {
      $\mathbf{u}$ \& $\mathbf{tree}$ \& $\mathbf{is\_leaf}$ \& $\mathbf{E}$ \& $\boldsymbol{\Theta}$ \& $\mathbf{N}^{\leftarrow}$ \& $\mathbf{N}^{\rightarrow}$ \& $\mathbf{score}$ \\
      $0$ \& $0$ \& $0$ \& $(0,0)$ \& $0$ \& $(0,0)$ \& $(0,0)$ \& $0$ \\
      $1$ \& $1$ \& $0$ \& $(1,0)$ \& $3$ \& $(0,0)$ \& $(B,B)$ \& $0$ \\
      $2$ \& $1$ \& $0$ \& $(0,1)$ \& $4$ \& $(0,0)$ \& $(3,B)$ \& $0$ \\
      $3$ \& $2$ \& $0$ \& $(1,0)$ \& $2$ \& $(0,0)$ \& $(B,B)$ \& $0$ \\
      $4$ \& $2$ \& $0$ \& $(0,1)$ \& $6$ \& $(2,0)$ \& $(B,B)$ \& $0$ \\
      $5$ \& $1$ \& $1$ \& $(0,0)$ \& $0$ \& $(0,0)$ \& $(3,4)$ \& $w_0$ \\
      $6$ \& $1$ \& $1$ \& $(0,0)$ \& $0$ \& $(0,4)$ \& $(3,B)$ \& $w_1$ \\
      |[fetchtxt]| $7$ \& |[fetchtxt]| $1$ \& |[fetchtxt]| $1$ \& |[fetchtxt]| $(0,0)$ \& |[fetchtxt]| $0$ \& |[fetchtxt]| $(3,0)$ \& |[fetchtxt]| $(B,B)$ \& |[fetchtxt]| $w_2$ \\
      $8$ \& $2$ \& $1$ \& $(0,0)$ \& $0$ \& $(0,0)$ \& $(2,B)$ \& $w'_0$ \\
      |[fetchtxt]| $9$ \& |[fetchtxt]| $2$ \& |[fetchtxt]| $1$ \& |[fetchtxt]| $(0,0)$ \& |[fetchtxt]| $0$
        \& |[fetchtxt]| $(2,0)$ \& |[fetchtxt]| $(B,6)$ \& |[fetchtxt]| $w'_1$ \\
      $10$ \& $2$ \& $1$ \& $(0,0)$ \& $0$ \& $(2,6)$ \& $(B,B)$ \& $w'_2$ \\
    };

    % Column background boxes
    \begin{scope}[on background layer]
      \foreach \c in {1,...,8}{
        \node[colbg, fit=(tbl-1-\c)(tbl-12-\c)] {};
      }
    \end{scope}

    % Horizontal separators: after header and after internal rows
    \foreach \r in {1,6}{
      \draw[gray!60, line width=.4pt]
        ([xshift=-2pt]tbl-\r-1.south west) -- ([xshift=2pt]tbl-\r-8.south east);
    }

    % Highlight fetched rows (one row per tree)
    \node[fetchrow, fit=(tbl-9-1)(tbl-9-8)] (fetchA) {};
    \node[fetchrow, fit=(tbl-11-1)(tbl-11-8)] (fetchB) {};

    % Query box
    \node[meta, anchor=north west, text width=3.35cm] (sample)
      at ([xshift=6mm]tbl.north east)
      {query $\mathbf{x}=(4,2)$\\[1mm]
       fetch one row per tree:\\
        $u_1=7,\;u_2=9$\\[1mm]
       and prove, for any fetched row $u$,\\
       $\mathbf{N}^{\leftarrow}[u]\le \mathbf{x}<\mathbf{N}^{\rightarrow}[u]$};

    % Arrows to the two fetched rows
    \draw[->, very thick, red!75!black]
      ([yshift=7pt]sample.west) to[out=180,in=0] ([xshift=2pt]fetchA.east);
    \draw[->, very thick, red!75!black]
      ([yshift=-7pt]sample.west) to[out=180,in=0] ([xshift=2pt]fetchB.east);

    % Outer frame
    \node[
      outer,
      fit=(ltitle)(rtitle)
          (t0r)(t0l)(t0rleaf)(t0w0)(t0w1)
          (t1r)(t1lleaf)(t1rr)(t1w1)(t1w2)
          (tbl)(sample)
    ] {};

  \end{tikzpicture}%
  }
  \caption{An illustration of the forest encoding and inference semantics.
  The left panel shows two compact trees with
  global node identifiers, and the right panel shows the committed
  objects used by inference. The highlighted entries are the fetched rows for one
  query, one per tree. Batched inference performs this fetch for every
  tree-sample pair in one shot and then aggregates scores across trees for
  each sample-class pair.}
  \Description{Two example decision trees and their committed table
  representation. The table includes node identifiers, tree identifiers,
  leaf indicators, split-feature masks, thresholds, hyper-box bounds, and leaf
  scores. Two leaf rows are highlighted to show how inference fetches one row
  per tree and checks interval containment for a query point.}
  \label{fig:forest-encoding-inference}
\end{figure*}
